\section{Economic Harm, Legal Scope, and Limits}
\label{sec:price_scope_conclusion}

The enforcement value of the screen does not depend on proving a
price effect. A screen can be useful if it ranks where forensic
attention should go, even when the price record is too coarsened
to support damages or causal incidence. Still, a cartel-adjacent
screen should leave some trace in the economic outcome agencies
care about. In BEC, frequent-loser presence is associated with
higher negotiated unit prices in the broad sample, but the
association changes sign under overlap restrictions. That tension
is informative. It says the price evidence should corroborate
the screening interpretation, not carry the main claim.

\subsection{Price Imprint}
\label{sec:price_imprint_v18}

In the broad item-level price regressions, frequent-loser
presence is positively associated with log negotiated unit price.
The coefficient is $\valTabPricesColOne$ with item and year fixed
effects, $\valTabPricesColTwo$ after adding procuring-unit fixed
effects, $\valTabPricesColThree$ in pregão, and
$\valTabPricesColFour$ in convite. Across the main descriptive
estimators, the implied range is $\valHeadlineRange$. The
pregão--convite split matters for interpretation. If frequent
losers were mainly quorum fillers created by the three-bidder
rule in convite, the price imprint should be stronger there.
Instead, the larger estimate appears in pregão, the real-time
electronic environment where loser-side positioning is easier to
deploy repeatedly.

\begin{table}[htbp]
\centering
\caption{Price Imprint and Scope}
\label{tab:price_scope_v18}
\footnotesize
\begin{adjustbox}{max width=\textwidth,center}
\begin{threeparttable}
\begin{tabular}{p{5.2cm}ccc}
\toprule
Exercise & Estimate & Standard error / $p$ & Interpretation \\
\midrule
Broad sample, item and year FE & $\valTabPricesColOne$ &
$\valTabPricesSEOne$ & Positive price imprint \\
Broad sample, item, year, PBU FE & $\valTabPricesColTwo$ &
$\valTabPricesSETwo$ & Positive after buyer controls \\
Pregão subsample & $\valTabPricesColThree$ &
$\valTabPricesSEThree$ & Larger in electronic auctions \\
Convite subsample & $\valTabPricesColFour$ &
$\valTabPricesSEFour$ & Smaller under quorum-prone procedure \\
\midrule
Overlap-cell ATT & $\valScopeOverlapCoef$ &
$p\valScopeOverlapPlt$ & Aggregate sign reversal \\
PS-trimmed ATT & $\valScopePSCoef$ &
--- & Stronger negative under trimming \\
Tender-value Q4, overlap ATT & $\valBetaQfour$ &
$p=\valBetaQfourP$ & Positive in high-value segment \\
Direct-CADE overlap items & $\valBetaCADEDirect$ &
$p=\valBetaCADEDirectP$ & Statistically null \\
\bottomrule
\end{tabular}
\begin{tablenotes}
\footnotesize
\item \textit{Notes:} Dependent variable is log negotiated unit
price. The broad-sample rows use the main BEC item-level sample.
The overlap rows restrict or reweight comparisons to cells with
treated and untreated support. The estimates are descriptive and
are not used as damages calculations.
\end{tablenotes}
\end{threeparttable}
\end{adjustbox}
\end{table}

The positive broad-sample association is not a causal estimate of
cover bidding's effect on prices. The available quasi-experimental
routes do not support that claim: discontinuity designs around
procedure thresholds lack the required density behavior, and the
policy-change design does not supply credible parallel trends.
For the argument here, the price evidence has a smaller job. It
asks whether the places highlighted by the screen also look
economically relevant. On the broad record they do.

\subsection{Where the Price Claim Stops}
\label{sec:price_scope_v18}

The overlap results draw the boundary. When comparisons are
restricted to cells where frequent-loser-present and
frequent-loser-absent items overlap on observables, the aggregate
coefficient reverses to $\valScopeOverlapCoef$; the PS-trimmed
estimate is $\valScopePSCoef$. Only $\valDroppedItemsShare$ of
treated items lack a within-cell counterfactual, so the reversal
is not mainly a dropped-observation story. It is a reweighting
story: the broad estimate and the overlap estimate answer
different questions about different parts of the item-value
distribution.

The segment decomposition makes that point explicit. In the
lowest three tender-value quartiles, the overlap estimates are
negative. In the highest tender-value quartile, the estimate is
positive: $\valBetaQfour$ ($p=\valBetaQfourP$). That is the
segment where cartel rents are most plausible ex ante because
the same fixed cost of coordination is spread over larger
contracts. The same table also shows a null estimate inside
direct-CADE overlap items, $\valBetaCADEDirect$
($p=\valBetaCADEDirectP$), which is another warning against
reading the price regressions as a clean treatment effect.

The mixed price record disciplines the legal claim. The evidence
says that the award-layer screen ranks cartel-adjacent
environments; it does not say that every flagged item is
overpriced, that the observed coefficient is a markup, or that
the screen identifies a damages base. Those are different legal
tasks and require different evidence.

\subsection{Portability and Strategic Response}
\label{sec:portability_adaptation_v18}

The screen is portable only where the institutional primitives
exist. A candidate jurisdiction needs public or auditable award
records with participant identities, winner identities, item
codes, and procurement timing. It also needs a procurement
environment in which persistent losing is economically
meaningful rather than purely mechanical. BEC fits that profile
particularly well: standardized goods are common, electronic
records are available, and pregão-style auctions preserve the
loser-side interaction environment in which the signal is
strongest. Settings dominated by bespoke services,
infrastructure, or sparse one-off procurement may still be
screenable, but the item-level price comparison and the
frequent-loser threshold would need to be rebuilt for the new
setting.

Strategic adaptation is the other limit. The screen is fairly
stable to simple rotation of cover bidders
($\valAUCAdaptRotation$ versus baseline $\valAUCBase$) and to
occasional wins that evade a strict zero-win rule
($\valAUCAdaptOccWins$). It weakens under CNPJ splitting
($\valAUCAdaptCNPJsplit$) and degrades more under
threshold-aware capping ($\valAUCAdaptThreshCap$). Combining
adaptations lowers AUC to $\valAUCAdaptCombined$. That pattern is
not surprising. A visible threshold invites threshold gaming. In
practice, the rule should be recalibrated, monitored for bunching
near the cutoff, and paired with bid-layer forensics where those
data can be recovered.

\subsection{Conclusion}
\label{sec:conclusion_v18}

Cartel enforcement often begins before the agency has the
evidence it would need to win a case. Procurement records then
pose a practical question: can the cheap layer do any legally
useful work before the expensive layer is opened? In São Paulo's
BEC platform, the answer is yes, but only within a disciplined
scope. Persistent zero-win participation ranks firms that are
cartel-adjacent on the loser side. It predicts adjudicated
cobidder exposure out of time, lines up with an economic
cover-bidder signature, and helps decide where bid-distribution
forensics should be run.

The contribution is not a shortcut to liability. It is a way to
order suspicion under incomplete information. That distinction is
why the screen belongs in the enforcement architecture rather
than in the proof stage. Used alone, it gives an agency a cheap
triage list. Used before bid-layer forensics, it reduces the
microdata footprint while preserving most of the adjudication
signal that a full joint model would recover. Used beyond that,
it would be overclaimed. The statistic identifies where legal
attention should start, not where legal responsibility ends.
